
\begin{table*}[th]
\centering
% \small{
\begin{tabular}{lccccccccc}
\toprule
\multirow{2}{*}{\textbf{Model}}  & \multicolumn{3}{c}{\textbf{SNLI}} & \multicolumn{3}{c}{\textbf{MultiNLI Matched}} &  \multicolumn{3}{c}{\textbf{MultiNLI Mismatched}}   \\
 & \textit{Full} & \hard & \easy & \textit{Full} & \hard & \easy & \textit{Full} & \hard & \easy  \\
\midrule\midrule
DAM  & 84.7 & 69.4 & 92.4 & 72.0 & 55.8 & 85.3 & 72.1 & 56.2 & 85.7 \\
ESIM & 85.8 & 71.3 & 92.6 & 74.1 & 59.3 & 86.2 & 73.1 & 58.9 & 85.2 \\ 
DIIN & 86.5 & 72.7 & 93.4 & 77.0 & 64.1 & 87.6 & 76.5 & 64.4 & 86.8  \\ 
\bottomrule
\end{tabular}
% }
\caption{Performance of high-performing NLI models on the full, \hard, and \easy NLI\com{ evaluation} test sets.
}
\label{tab:sotaresults}
\end{table*}


\section{Re-evaluating NLI Models}
\label{ssec:sota}

In Section~\ref{ssec:artifacts}, we showed that a model with no access to the premise can correctly classify many examples in both SNLI and MultiNLI, performing well above the most-frequent-class baseline.
This raises an important question about state-of-the-art NLI models: to what extent are they ``gaming'' the task by learning to detect annotation artifacts?

To answer this question, we partition each NLI test set into two subsets: examples that the premise-oblivious model classified accurately are labeled \easy, and those it could not are \hard. 


We then train an NLI model on the original training sets (from either SNLI or MultiNLI),\footnote{The MultiNLI models were trained on MultiNLI data alone (as opposed to a blend of MultiNLI and SNLI data).} and evaluate on the full test set, the \hard test set, and the \easy test set.
We ran this experiment on three high-performing NLI models:
the Decomposable Attention Model (DAM; \citealp{Parikh:2016}),\footnote{%Implementation from \texttt{AllenNLP}: \url{https://github.com/allenai/allennlp}
\url{github.com/allenai/allennlp}
} the Enhanced Sequential Inference Model (ESIM;  \citealp{Chen:2017}),\footnote{%Implementation from \citet{Williams:2017}:  \url{https://github.com/nyu-mll/multiNLI}
\url{github.com/nyu-mll/multiNLI}
} and the Densely Interactive Inference Network (DIIN; \citealp{Gong:2017}).\footnote{
% \url{github.com/YichenGong/Densely-Interactive-Inference-Network}}
\url{goo.gl/kCeZXm}} 
All models were retrained out of the box.



Table~\ref{tab:sotaresults} shows the performance of each model on the different splits.
While the models correctly classify some \hard examples, the bulk of their success is attributed to the \easy examples. 
This result implies that the ability of NLI models to recognize textual entailment is lower than previously perceived, and that such models rely heavily on annotation artifacts in the hypothesis to make their predictions.

\resolved{\omer{Maybe add later: An additional observation from this experiment is that DIIN performs significantly better than the other two models on \hard.}\roy{True, but it also performs better than the others on the full dataset, so this is not that surprising}}

A natural question to ask is whether it is possible to select a set of NLI training and test samples which do not contain easy-to-exploit artifacts.
One  solution might be to filter  \easy examples from the training set, retaining only  \hard examples.
However, initial experiments suggest that it might not be as straightforward to eliminate annotation artifacts once the dataset has been collected.

First, after removing the \easy examples, \hard examples might not necessarily be artifact-free.
For instance, removing all contradicting samples containing the word ``no'' (a strong indicator for contradiction, see \secref{analysis}), leaves the \hard dataset with this word mostly appearing in the neutral and entailing classes,
thus creating a new artifact.
Secondly, \easy examples contain important inference phenomena (e.g. the word ``animal'' is indeed a hypernym of ``dog''), and removing these examples may hinder the model from learning such phenomena.
Importantly, artifacts do not render any particular example {\it incorrect}; they are a problem with the sample distribution, which is skewed toward certain kinds of entailment, contradiction, and neutral hypotheses.\resolved{\nascomment{added this sentence, please check}}
Therefore, a better solution might not eliminate the artifacts altogether, but rather balance them across labels.
Future strategies for reducing annotation artifacts might involve experimenting with the prompts or training given to crowd workers, e.g., to encourage a wide range of strategies, or incorporating baseline or adversarial systems that flag examples that appear to use over-represented heuristics.
We defer research on hard-to-exploit NLI datasets to future work. 
